% Load data directly from JSON
\JSONParseFromFile[validate numbers=false,skip structures,externalize]{\fullData}{Data/full.json}
\JSONParseFromFile[validate numbers=false,skip structures,externalize]{\rgswData}{Data/rgsw.json}


\chapter{Results and Discussion}
\label{chp:results}
%%
%% INTRODUCTION
% ...
\vfill
I begin by examining 
\newpage

%%
%% PARAMETERS
\section{Parameters}
As a baseline I target a security parameter of $\lambda \approx 128$ 

\begin{table}[H]
    \caption{Base parameters for the }
    \label{tab:bench:params}
    \centering
    \begin{tabular}{lc}
        \toprule
        \textbf{Parameter}               & \textbf{Value}   \\
        \midrule
        decomposition $(\ell, B)$        & $(2, 2^{x})$ \\
        number of RNS moduli $k$         & $2$              \\
        standard deviation $\sigma$      & $3.19$           \\
        ciphertext modulus order $|Q|$   & $2^{120}$        \\
        plaintext modulus $t$            & $256$            \\
        polynomial degree $N$            & $4096$           \\
        security parameter $\lambda$     & $113.6$          \\
        \bottomrule
    \end{tabular}
\end{table}

% \begin{table}[!htb]
%     \caption{Base parameters for the small and large parameter sets}
%     \label{tab:bench:params}
%     \centering
%     \begin{tabular}{lcc}
%         \toprule
%         \textbf{Parameter} & \textbf{Small} & \textbf{Large} \\
%         \midrule
%         standard deviation $\sigma$      & $3.19$    & $3.19$ \\
%         number of RNS moduli $k$         & $2$       & $2$ \\
%         ciphertext modulus order $|Q|$   & $2^{120}$ & $2^{120}$ \\
%         plaintext modulus $t$            & $256$     & $65537$ \\
%         \hline % BV
%         decomposition $(\ell, B)$        & $(2, 2^{30}?)$  & $(2, ??)$ \\
%         polynomial degree $N$            & $4096$     & $16384$ \\
%         security parameter $\lambda$     & $113.6$    & $\approx 400$ \\
%         \hline % Hybrid
%         decomposition $(d_{num}, |P|)$   & $(1,120)$  & $(1, 120)$ \\
%         polynomial degree $N$            & $8192$     & $16384$ \\
%         security parameter $\lambda$     & $113.1$    & $\approx 240$ \\
%         \bottomrule
%     \end{tabular}
% \end{table}

\begin{itemize}
    \item For BV, the base $B = \lfloor \max q_i \rfloor$ is chosen automatically so to be the smallest $B$ that covers the largest RNS prime with $\ell$ digits when the program starts, because the $q_i$ are randomly chosen. 
    \item The decomposition parameter is $\ell = 2$, which was the smallest $\ell$ that could handle the noise.
\end{itemize}


\section{Noise Analysis}
\begin{itemize}
    \item How many users can be supported with these parameters?
    \item $d_{num}$ can probably be greater than 1, but this is not supported by the code yet (GHS is hard-coded at the RGSW level)
\end{itemize}

\subsubsection{External Product}
\begin{itemize}
    \item Take a fresh RGSW and apply it to an RLWE, save the result as the new RLWE and repeat
    \item Key takeaway: Hybrid performs better than BV at same security settings
    \item \textcolor{red}{Maybe a pointless section}
    % \item Noise becomes additive; grows linearly, as expected
    % \item At the base parameter set, unbounded but stays below 60 bits for at least 1000 iterations
\end{itemize}

% \definecolor{gaussblue}{RGB}{31,119,180}
% \definecolor{hpsgreen}{RGB}{44,160,44}
% \definecolor{canonorange}{RGB}{255,165,0}
% \definecolor{maxred}{RGB}{255,36,0}
% \definecolor{hpspurple}{RGB}{148,103,189}

% ---------- Compute fits in linear space ----------
% \pgfplotstableread[col sep=comma]{Data/Noise/ExternalProd/bv_small.csv}\bvsmall
% \pgfplotstablecreatecol[linear regression={x=n, y=noise}]{reg}{\bvsmall}
% \edef\smallSlope{\pgfplotstableregressiona}
% \edef\smallIntercept{\pgfplotstableregressionb}

\pgfplotstableread[col sep=comma]{Data/Noise/ExternalProd/bv.csv}\bvlarge
\pgfplotstablecreatecol[linear regression={x=n, y=noise}]{reg}{\bvlarge}
\edef\largeSlope{\pgfplotstableregressiona}
\edef\largeIntercept{\pgfplotstableregressionb}

\begin{figure}[H]
    \centering
    \caption{Accumulated noise from chained external products}
    \label{fig:bv:plot-1}
    \begin{tikzpicture}
    \begin{semilogyaxis}[
        noiseaxis,
        log basis y=2,
        xmin=0, xmax=1000,
        xlabel={Chain length},
        ylabel={$\log\|\epsilon\|_\infty$},
        legend pos=north west,
        set layers,
        width=\textwidth,
        height=0.55\textwidth,
    ]
        % ---------- Scatter plots ----------
        % \addplot[only marks, mark=*, mark size=1pt, color=gaussblue]
        %     table[col sep=comma, x=n, y=noise] {Data/Noise/ExternalProd/bv_small.csv};
        \addplot[only marks, mark=*, mark size=1pt, color=hpsgreen]
            table[col sep=comma, x=n, y=noise] {Data/Noise/ExternalProd/bv.csv};
            
        % ---------- Regression lines ----------
        % \addplot[thick, solid, dotted, on layer=axis foreground, color=canonorange, domain=1:1000, samples=300, forget plot]
        %     {\smallSlope*x + \smallIntercept};
        \addplot[thick, solid, dotted, on layer=axis foreground, color=hpspurple, domain=1:1000, samples=300, forget plot]
            {\largeSlope*x + \largeIntercept};
            
        % \legend{Small Parameters, Large Parameters}
        \legend{BV-RNS}
    \end{semilogyaxis}
    \end{tikzpicture}
\end{figure}


% \newpage

% \begin{tabular}{lc} 
%     \toprule
%     \textbf{n} & \textbf{noise} \\ % & \textbf{Role} \\ % Manual Table Header
%     \midrule
    
%     % Reads CSV, maps column headers to macros, and loops through rows
%     \csvreader[head to column names]{Data/Noise/ExternalProd/bv.csv}{}
%     {\n & \noise \\} 
    
%     % \bottomrule
% \end{tabular}



% \definecolor{gaussblue}{RGB}{31,119,180}
% \definecolor{hpsgreen}{RGB}{44,160,44}
% \definecolor{canonorange}{RGB}{255,165,0}
% \definecolor{maxred}{RGB}{255,36,0}
% \definecolor{hpspurple}{RGB}{148,103,189}

% ---------- Compute fits in linear space ----------
% \pgfplotstableread[col sep=comma]{Data/Noise/ExternalProd/hybrid_small.csv}\bvsmall
% \pgfplotstablecreatecol[linear regression={x=n, y=noise}]{reg}{\bvsmall}
% \edef\smallSlope{\pgfplotstableregressiona}
% \edef\smallIntercept{\pgfplotstableregressionb}

\pgfplotstableread[col sep=comma]{Data/Noise/ExternalProd/hybrid.csv}\bvlarge
\pgfplotstablecreatecol[linear regression={x=n, y=noise}]{reg}{\bvlarge}
\edef\largeSlope{\pgfplotstableregressiona}
\edef\largeIntercept{\pgfplotstableregressionb}

%% TODO: Ensure colors are consistent across graphs: BV ell = 2 should always be gaussblue etc.
%% TODO: Rename color names to be more memorizable
\begin{figure}[H]
    \centering
    \caption{Noise growth in chained external products.}
    \label{fig:hybrid:plot-1}
    \begin{tikzpicture}
    \begin{semilogyaxis}[
        noiseaxis,
        log basis y=2,
        xmin=0, xmax=300,
        xlabel={Chain length},
        ylabel={Error $\|\epsilon\|_\infty$},
        legend pos=south east,
        set layers,
        width=\textwidth,
        height=0.5\textwidth,
    ]
        % ---------- Scatter plots ----------
        \addplot[only marks, mark=*, mark size=.5pt, color=gaussblue]
            table[col sep=comma, x=n, y=noise] {Data/Noise/ExternalProd/bv_2.csv};
        \addplot[only marks, mark=*, mark size=.5pt, color=hpsgreen]
            table[col sep=comma, x=n, y=noise] {Data/Noise/ExternalProd/bv_3.csv};
        \addplot[only marks, mark=*, mark size=.5pt, color=canonorange]
            table[col sep=comma, x=n, y=noise] {Data/Noise/ExternalProd/bv_4.csv};
        \addplot[only marks, mark=*, mark size=.5pt, color=maxred]
            table[col sep=comma, x=n, y=noise] {Data/Noise/ExternalProd/bv_5.csv};
        \addplot[only marks, mark=*, mark size=.5pt, color=hpspurple]
            table[col sep=comma, x=n, y=noise] {Data/Noise/ExternalProd/hybrid.csv};
            
        \legend{BV $\ell = 2$, BV $\ell = 3$, BV $\ell = 4$, BV $\ell = 5$, Hybrid}
    \end{semilogyaxis}
    \end{tikzpicture}
\end{figure}

\subsubsection{Internal product}
\begin{itemize}
    \item Hybrid supports at most 1 single internal product before the noise blows up
\end{itemize}

\definecolor{plotblue}{HTML}{00ADB5}
% \definecolor{hpsgreen}{RGB}{44,160,44}
% \definecolor{canonorange}{RGB}{255,165,0}
% \definecolor{maxred}{RGB}{255,36,0}
% \definecolor{hpspurple}{RGB}{148,103,189}

% ---------- Compute fits in linear space ----------
% \pgfplotstableread[col sep=comma]{Data/Noise/InternalProd/bv_small.csv}\bvsmall
% \pgfplotstablecreatecol[linear regression={x=n, y=noise}]{reg}{\bvsmall}
% \edef\smallSlope{\pgfplotstableregressiona}
% \edef\smallIntercept{\pgfplotstableregressionb}

\pgfplotstableread[col sep=comma]{Data/Noise/InternalProd/bv_2.csv}\bvlarge
\pgfplotstablecreatecol[linear regression={x=n, y=noise}]{reg}{\bvlarge}
\edef\intSlope{\pgfplotstableregressiona}
\edef\intIntercept{\pgfplotstableregressionb}

\begin{figure}[H]
    \centering
    \caption{Accumulated noise from chained internal products.}
    \label{fig:bv:plot-2}
    \begin{tikzpicture}
    \begin{semilogyaxis}[
        noiseaxis,
        log basis y=2,
        xmin=0, xmax=1000,
        xlabel={Chain length},
        ylabel={Error $\|\epsilon\|_\infty$},
        legend pos=south east,
        set layers,
        width=\textwidth,
        height=0.6\textwidth,
    ]
        % ---------- Scatter plots ----------
        % \addplot[only marks, mark=*, mark size=1pt, color=gaussblue]
        %     table[col sep=comma, x=n, y=noise] {Data/Noise/InternalProd/bv_small.csv};
        % \addplot[only marks, mark=*, mark size=1pt, color=gaussblue]
        %     table[col sep=comma, x=n, y=noise] {Data/Noise/InternalProd/bv_1.csv};
        \addplot[only marks, mark=*, mark size=1pt, color=hpsgreen]
            table[col sep=comma, x=n, y=noise] {Data/Noise/InternalProd/bv_2.csv};
        \addplot[only marks, mark=*, mark size=1pt, color=canonorange]
            table[col sep=comma, x=n, y=noise] {Data/Noise/InternalProd/bv_3.csv};
        \addplot[only marks, mark=*, mark size=1pt, color=gaussblue]
            table[col sep=comma, x=n, y=noise] {Data/Noise/InternalProd/bv_4.csv};
            
        % ---------- Regression lines ----------
        % \addplot[thick, solid, dotted, on layer=axis foreground, color=canonorange, domain=1:1000, samples=300, forget plot]
        %     {\smallSlope*x + \smallIntercept};
        % \addplot[thick, solid, dotted, on layer=axis foreground, color=hpspurple, domain=1:1000, samples=300, forget plot]
        %     {\largeSlope*x + \largeIntercept};
        
        % ---------- Decryption failure ---------
        % \addplot[color=maxred, domain=1:1000]{2^120};
            
        \legend{$\ell = 2$ (baseline), $\ell = 3$, $\ell = 4$}
    \end{semilogyaxis}
    \end{tikzpicture}
\end{figure}


\subsubsection{sPAR}
\begin{itemize}
    \item If the parameters support even the base case of $\eta = 2$ users, then the scheme can support a nearly arbitrary amount of users with the same parameters
    \item Noise is a big problem due to the internal product
    \item Getting enough noise headroom at $\lambda = 128$ requires very large parameters, making the protocol much slower than anticipated
\end{itemize}


%%
%% BENCHMARKS
\section{Benchmarks}
\begin{itemize}
    \item Run on AMD Ryzen 5 5600X CPU 6-core/12-thread@3.7 GHz with 32Gb RAM in WSL2 Ubuntu Linux 
    \item Compare numbers with parity in the security parameter (according to the \href{https://github.com/malb/lattice-estimator#}{Lattice Estimator \cite{cryptoeprint:2015/046}})
    \item Compare numbers at SIMD-capable level
\end{itemize}

%%
%% Encrypt RGSW, external product, internal product
\subsection{TFHE Operations}
\begin{table}[!htb]
    \begin{subtable}{.5\linewidth}
        \begin{tabular}{cc}
        \textbf{Operation} & \textbf{Time} \\
            Encrypt RGSW        & \num{\jsonms{\rgswData}{benchmarks[0].cpu_time}}~ms \\
            External Product    & \num{\jsonms{\rgswData}{benchmarks[1].cpu_time}}~ms \\
            Internal Product    & \num{\jsonms{\rgswData}{benchmarks[2].cpu_time}}~ms
        \end{tabular}
        \caption{BV}
        \label{tab:bench:operations:bv}
    \end{subtable}%
    \begin{subtable}{.5\linewidth}
      \centering
        \begin{tabular}{cc}
        \textbf{Operation} & \textbf{Time} \\
            Encrypt RGSW        & \num{\jsonms{\rgswData}{benchmarks[3].cpu_time}}~ms \\
            External Product    & \num{\jsonms{\rgswData}{benchmarks[4].cpu_time}}~ms \\
            Internal Product    & \num{\jsonms{\rgswData}{benchmarks[5].cpu_time}}~ms\footnote{Not really valid since the noise is too large/result is incorrect...}
        \end{tabular}
        \caption{GHS/Hybrid}
        \label{tab:bench:operations:ghs}
    \end{subtable} 
    \caption{RGSW/TFHE Operations}
\end{table}

\subsubsection{Multi-Threaded Performance}
\begin{itemize}
    \item Works now, but only some functions (EncryptRGSW, BV Decompose) are optimized to take advantage of multi-threading.
    \item $\leq k\ell$ threads saturate, so $k = 2, \ell = 3$ has optimal performance at 6 threads.
    \item Ideally more threads would used to perform multiple coefficients simultaneously, but this is not implemented.
\end{itemize}

\textbf{Table of improvements +X\% at 1 to 6 threads}

\subsubsection{sPAR}
%%
%% Client::Encrypt, Server::Write, Client::Decrypt, Server::Decrypt
\begin{table}[H]
    \caption{One round of the sPAR protocol, \textcolor{red}{using the \textbf{outdated} GHS gadget}}
    \label{tab:bench:full}
    \centering
    \begin{tabular}{cccccc}
        \toprule
        $\mathbf{\eta}$ & \textbf{Total time} & \textbf{Encrypt} & \textbf{Write} & \textbf{Partial Dec.} & \textbf{Server Dec.} \\
        \midrule
        \JSONParseArrayMapInline{\fullData}{benchmarks}{
          $\nvalue{#1}$ &
          \num{\jsonms{\fullData}{benchmarks[#1].cpu_time}} & 
          \num{\jsonmspern{\fullData}{benchmarks[#1].encrypt}{\nvalue{#1}}} &
          \num{\jsonmspern{\fullData}{benchmarks[#1].write}{\nvalue{#1}}} &
          \num{\jsonmspern{\fullData}{benchmarks[#1].partial}{\nvalue{#1}}} &
          \num{\jsonms{\fullData}{benchmarks[#1].final}} \\
        }
        \bottomrule
    \end{tabular}
\end{table}



\section{Discussions}
The results show that 
